% PAGES: 284-284
% Self-contained Example:/Answer: pair opens and closes within this single
% page (folio 268); the closing square is printed at the very end of this
% file, after "...largest expected return of any of the four assets."

The maximum expected return of a portfolio with variance of return equal to
$\sigma_P^2$ is
\begin{align}
\mu_{max} &= r_f+\bar\mu^tw_{max} \;=\;
\sigma_P\frac{\bar\mu^t\Sigma_R^{-1}\bar\mu}{\sqrt{\bar\mu^t\Sigma_R^{-1}\bar\mu}}
\notag\\
&= r_f+\sigma_P\sqrt{\bar\mu^t\Sigma_R^{-1}\bar\mu}.
\end{align}

The pseudocode for the asset allocation for the maximum return portfolio computed
using the formulas (9.59--9.61) can be found in Table 9.3.

\par\medskip\noindent\textit{Example:}\ Find the \$100 million maximum return
portfolio with 27\% standard deviation of return invested in cash and four assets
with expected returns equal to 2\%, 1.75\%, 2.5\%, and 1.5\%, and with the following
covariance matrix of their returns:
\[
\begin{pmatrix}
0.09 & 0.01 & 0.03 & -0.015 \\
0.01 & 0.0625 & -0.02 & -0.01 \\
0.03 & -0.02 & 0.1225 & 0.02 \\
-0.015 & -0.01 & 0.02 & 0.0576
\end{pmatrix}.
\]

Assume that the risk free rate is 1\%.

\begin{answerx}
We know that $\mu_1=0.02$, $\mu_2=0.0175$, $\mu_3=0.025$, $\mu_4=0.015$, and $r_f=
0.01$. We look for the asset and cash allocation corresponding to the maximum return
portfolio with standard deviation of return equal to $\sigma_P = 0.27$.

Note that $\bar\mu=\mu-r_f\bfone$ and $\Sigma_R$ are given by
\[
\bar\mu \;=\; \begin{pmatrix} 0.01 \\ 0.0075 \\ 0.015 \\ 0.005 \end{pmatrix}; \quad
\Sigma_R \;=\; \begin{pmatrix}
0.09 & 0.01 & 0.03 & -0.015 \\
0.01 & 0.0625 & -0.02 & -0.01 \\
0.03 & -0.02 & 0.1225 & 0.02 \\
-0.015 & -0.01 & 0.02 & 0.0576
\end{pmatrix}.
\]

The vector $\Sigma_R^{-1}\bar\mu$ is computed by using the Cholesky decomposition to
solve the linear system corresponding to the matrix $\Sigma_R$ and to the right hand
side $\bar\mu$, i.e., $\Sigma_R^{-1}\bar\mu =
\text{linear\_solve\_cholesky}(\Sigma_R,\bar\mu)$; see Table 6.2 and also the
pseudocode from Table 9.3. From (9.59) and (9.60), we find that the asset weights
vector corresponding to the maximum return portfolio is
\[
w_{max} \;=\;
\frac{\sigma_P}{\sqrt{\bar\mu^t\Sigma_R^{-1}\bar\mu}}\Sigma_R^{-1}\bar\mu \;=\;
\begin{pmatrix} 0.2938 \\ 0.6773 \\ 0.4905 \\ 0.3891 \end{pmatrix},
\]
and the weight of the cash position is
\[
w_{max,cash} \;=\; 1-\bfone^tw_{max} \;=\; -0.8507.
\]

Thus, the \$100 million maximum return portfolio with 27\% standard deviation of
return is obtained by borrowing \$85.07 million and investing \$29.38 million in
asset 1, \$67.73 million in asset 2, \$49.05 million in asset 3, and \$38.91 million
in asset 4.

From (9.61), we find that the maximum expected return of this portfolio is
$\mu_{max} = 0.0273$, i.e., 2.73\%, which is more than 2.5\%, the largest expected
return of any of the four assets. \hfill$\square$
\end{answerx}
